\documentclass[12pt]{article}
\usepackage[utf8]{inputenc}
\usepackage[spanish]{babel}
\usepackage{graphicx}
\usepackage{geometry}
\usepackage{pdfpages}
\geometry{margin=2.5cm}
\usepackage[utf8]{inputenc}
\usepackage{booktabs} % Para líneas de tabla más elegantes
\usepackage{tabularx} % Para que la tabla se ajuste al ancho del texto
\usepackage{makecell} % más de una línea en la misma celda dentro de una tabla

\title{\Huge Bitácora 2\\ \Large Análisis de la relación entre el uso de redes sociales y el rendimiento académico en estudiantes}
\author{Katia Moreno Aspiroz C6A060\\
Jeremy Garcia Solano C33171\\
Luis Diego Elizondo Fennell C4E844\\
Sebastián Calderón Segura C21517}
\date{\today}

\begin{document}
\maketitle
\newpage

\tableofcontents
\newpage

\section{Introducción}
\section{Planificación}
\subsection{Respuesta retroalimentación}
necesitamos la retroalimentación
\subsection{Rastro de decisiones}
\subsection{Audiencia del proyecto}
Esta semana, le diría a un profesor universitario que ...

\section{Ordenamiento Literatura}
\subsection{Fichas literarias}

\section*{Ficha 1}

\subsection*{Nivel descriptivo}
\begin{itemize}
  \item \textbf{Título:} Redes Sociales y rendimiento académico en estudiantes de Administración, Universidad Nacional de Trujillo. Perú
  \item \textbf{Autores:} Sandra Lizzette León Luyo, Denis Guizela Chávez Bejarano, Nelly Victoria De La Cruz Ruiz, Heyner Yuliano Marquez Yauri
  \item \textbf{Año:} 2023
  \item \textbf{Nombre del tema:} Relación entre redes sociales y rendimiento académico en estudiantes universitarios durante la pandemia.
  \item \textbf{Clasificación:} Metodológica (se centra en un estudio correlacional descriptivo con cuestionarios).
  \item \textbf{Resumen en una oración:} Estudia si el uso de redes sociales afecta el rendimiento académico en estudiantes de Finanzas.
  \item \textbf{Argumento central:} No existe relación positiva significativa entre redes sociales y rendimiento académico, aunque algunos estudiantes perciben beneficios.
  \item \textbf{Problemas:} Limitaciones de conectividad y tamaño de muestra reducida (46 estudiantes).
\end{itemize}

\subsection*{Nivel analítico}
\begin{itemize}
  \item \textbf{Conexión con nuestro proyecto:} Relaciona directamente redes sociales con rendimiento académico. Sin embargo, centra el estudio en una población reducida de estudiantes de una carrera universitaria en específico.
  \item \textbf{Lo que NO dice el autor:} No analiza diferencias por tipo de red social más allá de frecuencia ni efectos longitudinales.
  \item \textbf{Contraste con otra fuente:} A diferencia de Aljemely 2020, esta ficha no relaciona el uso de redes sociales con efectos negativos importantes en el rendimiento académico, sino que encuentra resultados más neutrales.
  \item \textbf{Evaluación de aplicabilidad (1-5):} 4 - metodología correlacional aplicable, aunque con muestra pequeña.
  \item \textbf{Pregunta al autor:} ¿Cómo cambiarían los resultados si se midiera rendimiento con indicadores más amplios que solo notas de la carrera de Finanzas?
  \item \textbf{Resumen argumentativo:} El artículo busca determinar si redes sociales afectan el rendimiento académico durante la pandemia. Aunque concluye que no hay relación significativa, sí identifica percepciones positivas en un 19.6\% de estudiantes. Para nuestro proyecto, aporta un marco metodológico de cuestionarios y variables de tiempo de uso y concentración, pero carece de análisis longitudinal. Se usará como evidencia de estudios que encuentran relaciones débiles o nulas.
\end{itemize}

\section*{Ficha 2}

\subsection*{Nivel descriptivo}
\begin{itemize}
  \item \textbf{Título:} Coorrelación del uso de redes sociales con el rendimiento académico de estudiantes de nivel superior.
  \item \textbf{Autor:} Jesús Adolfo Rodelo Moreno, Jorge Lizárraga Reyes
  \item \textbf{Año:} 2018
  \item \textbf{Nombre del tema:} Relación entre el uso de redes sociales y el rendimiento académico usando análisis de correlación.
  \item \textbf{Clasificación:} Metodológica (usa métodos de estadísticos de correlación a partir de encuestas).
  \item \textbf{Resumen en una oración:} Estudia si existe relación entre el uso de redes sociales y el rendimiento académico en estudiantes universitarios.
  \item \textbf{Argumento central:} El uso de redes sociales puede tener relación positiva con el rendimiento académico cuando se utilizan con fines académicos.
  \item \textbf{Problemas:} La muestra es limitada y los datos dependen de lo que responden los estudiantes en las encuestas.
\end{itemize}

\subsection*{Nivel analítico}
\begin{itemize}
  \item \textbf{Conexión con nuestro proyecto:} Se relaciona bastante con nuestro trabajo porque también usa encuestas y análisis de correlación para estudiar redes sociales y rendimiento académico.
  \item \textbf{Lo que NO dice el autor:} No analiza otros factores que podrían influir en las notas, como hábitos de estudio o situación personal de cada estudiante.
  \item \textbf{Contraste con otra fuente:} A diferencia de León Luyo 2023, este estudio sí encuentra algunas relaciones positivas entre redes sociales y rendimiento académico.
  \item \textbf{Evaluación de aplicabilidad (1-5):} 5 - muy útil para el proyecto porque utiliza una metodología parecida a la nuestra.
  \item \textbf{Pregunta al autor:} ¿Qué tipo de uso de redes sociales creen que influye más en el rendimiento académico, el académico de forma positiva o en recreativo de forma negativa?
  \item \textbf{Resumen argumentativo:} El artículo analiza la relación entre el uso de redes sociales y el rendimiento académico mediante encuestas y análisis de correlación. Es útil para este proyecto porque trabaja con variables muy similares a las nuestras y también utiliza una metodología similar. Además, nos sirve para comparar resultados con estudios distintos que encuentran relaciones diferentes entre redes sociales y desempeño académico.
\end{itemize}

\section*{Ficha 3}

\subsection*{Nivel descriptivo}
\begin{itemize}
  \item \textbf{Título:} Relación entre adicción a teléfonos inteligentes, redes sociales, alienación y rendimiento académico entre estudiantes de pregrado
  \item \textbf{Autor:} Yousef A. Aljemely
  \item \textbf{Año:} 2020
  \item \textbf{Nombre del tema:} Impacto de la adicción a smartphones y redes sociales en rendimiento académico.
  \item \textbf{Clasificación:} Temática (se centra en adicción tecnológica y sus efectos).
  \item \textbf{Resumen en una oración:} Analiza cómo la adicción a smartphones y redes sociales afecta alienación y rendimiento académico.
  \item \textbf{Argumento central:} La adicción a redes sociales se relaciona con alienación y puede disminuir rendimiento académico, aunque hay efectos positivos menores.
  \item \textbf{Problemas:} Resultados limitados a una muestra pequeña y contexto específico.
\end{itemize}

\subsection*{Nivel analítico}
\begin{itemize}
  \item \textbf{Conexión con nuestro proyecto:} Relevante porque introduce la variable “adicción” como mediadora entre redes sociales y rendimiento.
  \item \textbf{Lo que NO dice el autor:} No diferencia entre tipos de redes sociales ni actividades académicas específicas.
  \item \textbf{Contraste con otra fuente:} Contrasta con León Luyo 2023, que no encuentra relación significativa; aquí sí se reporta impacto negativo.
  \item \textbf{Evaluación de aplicabilidad (1-5):} 4 - metodología descriptiva aplicable, aunque con sesgo cultural.
  \item \textbf{Pregunta al autor:} ¿Cómo mediría la adicción de manera objetiva más allá de cuestionarios de percepción?
  \item \textbf{Resumen argumentativo:} El artículo aborda el problema de la adicción tecnológica y su impacto en rendimiento académico. Aunque encuentra efectos negativos, también señala un leve impacto positivo en GPA. Para nuestro proyecto, evidencia que estudios pueden encontrar una relación entre el uso de redes sociales y el rendimiento académico. A su vez, aporta la dimensión de “adicción” como variable que puede explicar diferencias en resultados contradictorios entre estudios.
\end{itemize}

\section*{Ficha 4}

\subsection*{Nivel descriptivo}
\begin{itemize}
  \item \textbf{Título:} Influence of Social Media on the Academic Performance Among Undergraduate Students of University of Ilorin.
  \item \textbf{Autor:} Olamilekan O. Ogunbiyi.
  \item \textbf{Año:} 2024
  \item \textbf{Nombre del tema:} Influencia del uso de redes sociales en el rendimiento académico de estudiantes universitarios.
  \item \textbf{Clasificación:} Metodológica (usa encuestas y análisis estadístico para estudiar la relación entre redes sociales y rendimiento académico).
  \item \textbf{Resumen en una oración:} Analiza cómo el uso de redes sociales influye en el rendimiento académico de estudiantes universitarios mediante encuestas y análisis estadístico.
  \item \textbf{Argumento central:} Las redes sociales pueden ser útiles para el aprendizaje, pero un uso excesivo puede afectar negativamente al rendimiento académico y la concentración.
  \item \textbf{Problemas:} El estudio se basa en respuestas subjetivas de estudiantes de una sola universidad.
\end{itemize}

\subsection*{Nivel analítico}
\begin{itemize}
  \item \textbf{Conexión con nuestro proyecto:} Se relaciona directamente con nuestro trabajo porque estudia variables similares como tiempo de uso y rendimiento académico.
  \item \textbf{Lo que NO dice el autor:} No explica si el tiempo de uso afecta igual a todos los estudiantes o si depende de cómo utilicen las redes sociales.
  \item \textbf{Contraste con otra fuente:} A diferencia de León Luyo 2023, este estudio sí relaciona el uso excesivo de redes sociales con una bajada en el rendimiento académico.
  \item \textbf{Evaluación de aplicabilidad (1-5):} 4 - útil para comparar resultados y metodología con nuestro proyecto.
  \item \textbf{Pregunta al autor:} ¿Los resultados serían diferentes si el estudio se realizase en estudiantes de otras carreras o países?
  \item \textbf{Resumen argumentativo:} El artículo estudia cómo el uso de redes sociales puede influir tanto positiva como negativamente en el rendimiento académico de estudiantes universitarios. Utiliza encuestas y análisis estadístico para relacionar tiempo de uso, concentración y desempeño académico. Para nuestro proyecto es útil porque cuenta con variables similares a las de nuestra base de datos y permite comparar resultados de forma más clara.
\end{itemize}


\subsection{Ordenamiento Literatura}
Con el objetivo de organizar ...

\begin{table}[h]
\centering
\begin{tabular}{|c|c|c|c|c|}
\hline
\multicolumn{2}{|c|}{Organización} & \multicolumn{3}{c|}{Literatura}\\
\cline{1-5}
Clasificación & Justificación breve & Título & Año & Autor(es) \\
\hline
Metodológica & \makecell{Estudia si el uso de\\ redes
sociales afecta el\\ rendimiento académico en\\ estudiantes de Finanzas} & \makecell{Redes Sociales y rendimiento\\ académico en estudiantes de\\ Administración, Universidad\\ Nacional de Trujillo. Perú} & 2003 & \makecell{Sandra L. L.,\\ Denis G. C.,\\ Nelly V. C.,\\ Heyner Y. M.}\\
Teórica & \makecell{Examina cómo la\\ estructura y contenido\\ de redes de amistad\\ afectan el rendimiento\\ académico} & \makecell{Análisis de redes sociales y\\ rendimiento académico:\\ lecciones a partir del caso\\ de los Estados Unidos} & 2010 & Martín Santos\\
Temática & \makecell{Analiza cómo la adicción\\ a smartphones y redes\\ sociales afecta alienación y\\ rendimiento académico} & \makecell{Relación entre adicción a\\ teléfonos inteligentes, redes\\ sociales, alienación\\ y rendimiento académico\\ entre estudiantes de\\ pregrado} & 2020 & \makecell{Yousef A.\\ Aljemely}\\
\hline
\end{tabular}
\caption{Ordenamiento de la literatura.}
\end{table}

\section{Enlaces de la literatura}
cada párrafo con resumen, contraste y contribución propia.

\section{Análisis estadísticos}
\subsection{Identidad visual}
\subsection{Análisis descriptivo}
\subsection{Fichas de resultados}
\subsection{Imagen que cuenta la historia}

\section{Arco narrativo}
al explorar los datos, descubrimos que...
\section{Diario de aprendizaje}
\section{Log de contribuciones}

\section{Registro de uso de IA}
\begin{itemize}
    \item Se utiliza la IA para crear el cuadro correspondiente al ordenamiento de la literatura con la estructura que se propone en la Bitácora 2:
    \begin{itemize}
        \item \textbf{Herramienta:} ChatGPT
        \item \textbf{Promt:} Necesito que me expliques cómo crear en lenguaje Latex una tabla que contenga celdas combinadas.
        \item \textbf{Respuesta:} "Puedes hacerlo combinando \texttt{multicolumn}, \texttt{cline} y líneas parciales La clave es:\\
        $|$ en la definición de columnas $\rightarrow$ líneas verticales.\\
        \texttt{hline} $\rightarrow$ línea horizontal completa.\\
        \texttt{cline\{i-j\}} $\rightarrow$ línea horizontal solo entre ciertas columnas.\\
        \texttt{multicolumn} $\rightarrow$ unir columnas y controlar si aparecen líneas alterales."\\
        Y brinda un ejemplo de código.
        \item \textbf{Evaluación:} Muy útil esta ayuda ya que la información de tablas de los manuales de Latex son básicas y no contienen información sobre la función \texttt{multicolumn}.
        \item \textbf{Modificaciones:} Se modificó la tabla con la información necesaria para el trabajo y las celdas combinadas que requeríamos en este caso, que no correspondían a las que combinó la IA.
        \item \textbf{Aprendizaje:} Realizar cuadros combinados en LaTeX.
    \end{itemize}
    
\end{itemize}


\end{document}
